\chapter{Chapter 6 Refresher: Pretraining, GPT-Style Models, and In-Context Learning --- Pretraining Objective and Decoding Controls}

\section*{Setting the Scene}
Training optimises Preface Steps~8--10 on teacher-forced contexts; deployment samples from the learned conditional family under budgets and safety envelopes---same \(\theta\), different \emph{pushforward} map from logits to tokens.

\paragraph{Step A --- Autoregressive likelihood reminder.}
\textbf{Problem.} Guarantee coherence across positions while fitting data.
\textbf{Insight.} Factorisation \(\prod_t p_\theta(x_t\mid x_{<t})\) ties optimisation to causal graphs already rehearsed in Preface Step~3.
\textbf{Lineage.} Shannon's entropy estimates for English $\to$ modern causal transformers.
\textbf{Mental model.} Teacher forcing mismatches free-running drift---monitor both regimes.

\paragraph{Step B --- Inference-time reshaping.}
\textbf{Problem.} Raw softmax may be too peaked or too flat for products.
\textbf{Insight.} Temperature, top-\(k\), top-\(p\) apply measure transforms \emph{after} forward pass---no gradient necessarily flows (pure inference policy).
\textbf{Lineage.} Nucleus sampling (Holtzman et al.) formalises mass truncation psychology.
\textbf{Mental model.} Treat decoding knobs as UI levers on Shannon entropy (Preface Step~9).

\paragraph{Step C --- Scaling budgets.}
\textbf{Problem.} Engineers trade parameters vs tokens under compute roofs.
\textbf{Insight.} Empirical iso-loss laws operationalise Preface Step~7 qualitatively at trillion-token scales.
\textbf{Lineage.} Kaplan (2020); Chinchilla (2022).
\textbf{Mental model.} Laws fit historical pipelines---extrapolate cautiously when data composition shifts.

Decoding is theatre lighting on the same actor (\(\theta\)); training chooses the script likelihood, decoding chooses spotlight intensity.

\section*{Notation Introduced in This Chapter}
\begin{itemize}[leftmargin=1.5em]
  \item $z_i$: logit for candidate token $i$.
  \item $p_i$: probability of candidate token $i$ after normalization.
  \item $\tau$: temperature parameter.
  \item $k$: top-$k$ truncation count.
  \item $p$: top-$p$ cumulative mass threshold.
  \item $N,D$: model/data scale variables in scaling discussions.
\end{itemize}

\section*{What You Need To Recall Precisely}
\begin{itemize}[leftmargin=1.5em]
  \item \textbf{Autoregressive NLL objective.}
  For corpus sequences \(x_{1:T}\), training maximises \(\sum_t \log p_\theta(x_t\mid x_{<t})\), equivalently minimises per-token cross-entropy averaged over positions and documents. This is the Chapter~1 chain rule factorisation evaluated inside the transformer stack that produces conditional softmax parameters at each \(t\).
  \item \textbf{Temperature transform:}
  \[
    p_i(\tau)=\frac{\exp(z_i/\tau)}{\sum_j \exp(z_j/\tau)}.
  \]
  \item \textbf{Top-$k$ and top-$p$ (ordered operations).}
  Start from full softmax masses \(p_i\) over vocabulary.
  \begin{enumerate}[leftmargin=1.5em]
    \item \textbf{Top-$k$:} keep the $k$ largest probabilities, zero the rest, divide remaining masses by their sum so they again form a distribution on the surviving support.
    \item \textbf{Top-$p$ (nucleus):} sort probabilities descending, greedily add tokens until cumulative mass reaches threshold $p\in(0,1]$, zero outside that prefix, renormalise on the kept set.
  \end{enumerate}
  Both policies intentionally discard probability mass outside a subset of vocabulary entries before sampling; they reshape diversity without changing trained weights.
  \item \textbf{Objective/metric distinction:} training objective may not align perfectly with deployment metric.
  \item \textbf{Scaling fits:} compute/data/parameter fronts are empirical and objective-dependent.
\end{itemize}

\section*{Pointers}
\begin{itemize}[leftmargin=1.5em]
  \item CE/KL/perplexity: see Chapter 1 refresher.
  \item Optimizer and scaling-law mechanics: see Chapter 2 refresher.
  \item Attention architecture assumptions: see Chapter 5 refresher.
\end{itemize}

\section*{How These Results Build on Each Other}
\begin{enumerate}[leftmargin=1.5em]
  \item Train model to minimize autoregressive NLL.
  \item Convert logits to probabilities at inference.
  \item Apply decoding transform (temperature/truncation).
  \item Sample outputs from transformed distribution.
  \item Evaluate behavior with metrics that may differ from training objective.
\end{enumerate}

\section*{Reminder Glossary (Term by Term)}
\begin{itemize}[leftmargin=1.5em]
  \item \textbf{Autoregressive objective:} predict each token conditioned on prefix.
  \item \textbf{Decoding policy:} rule for turning model distribution into emitted tokens.
  \item \textbf{Temperature:} entropy-control knob on logits.
  \item \textbf{Top-$k$:} fixed-cardinality truncation.
  \item \textbf{Top-$p$:} mass-based truncation.
  \item \textbf{Scaling law:} empirical relation between loss and resource axes.
\end{itemize}

\section*{Worked Mini Example (Temperature)}
Take logits $z=(2,1,0)$.
With $\tau=1$:
\[
p\approx(0.665,0.245,0.090).
\]
With $\tau=2$:
\[
p\approx(0.506,0.307,0.186).
\]
Higher temperature flattens the distribution, increasing diversity and risk simultaneously.

\section*{Worked Example (Decode Policy Choice for Same Model)}
Suppose next-token distribution after softmax is:
\[
(0.52,0.24,0.12,0.07,0.05).
\]
Under top-$k$ with $k=2$, support becomes first two tokens, renormalized to
\[
\left(\frac{0.52}{0.76},\frac{0.24}{0.76}\right)\approx(0.684,0.316).
\]
Under top-$p$ with $p=0.9$, first four tokens are retained (mass $0.95$), renormalized to:
\[
(0.547,0.253,0.126,0.074).
\]
Same trained model, different generation behavior due only to policy.

\section*{Intuition Checkpoints}
\begin{itemize}[leftmargin=1.5em]
  \item Decoding does not retrain the model; it samples from a transformed output distribution.
  \item Top-$k$ and top-$p$ can behave similarly in peaked distributions and very differently in flat distributions.
  \item Apparent ``emergence'' can be caused by metric discontinuities, not abrupt capability phase transitions.
\end{itemize}

\section*{Further Refresh}
\begin{itemize}[leftmargin=1.5em]
  \item \textbf{Hoffmann et al. (2022)} --- compute-optimal scaling interpretation.
  \item \textbf{Holtzman et al.} --- nucleus sampling and degeneration analysis.
  \item Standard decoding notes from CS224N/modern LLM systems courses.
\end{itemize}
